\section{System Architecture \& Constrained Engineering}
\label{sec:arch}

Fusion's architecture is designed around a single principle: \textbf{zero unnecessary abstraction, zero dynamic runtime dependency, and zero C bindings}.

\subsection{Pure Rust Runtime \& Android Bionic Native Compilation}
Standard agent harnesses rely on Python or Node.js runtimes. On Android, executing Node.js requires installing Termux and running under glibc compatibility wrappers or proot environments. This introduces three major bottlenecks:
\begin{enumerate}
    \item System call interception in proot adds $2.5\times\text{--}4\times$ execution overhead on filesystem operations (e.g., git status or file searches).
    \item OpenSSL dynamic link failures frequently occur because Android Bionic libc does not supply standard desktop certificates or symbol versions.
    \item Resident set size (RSS) of the V8 JavaScript engine or Python interpreter exceeds 200~MB before the agent even loads its system prompt.
\end{enumerate}

Fusion completely eliminates these bottlenecks by compiling to a single, static 6.8~MB native binary via the standard Rust Android NDK toolchain:
\begin{lstlisting}[language=bash]
cargo build --target aarch64-linux-android --release
\end{lstlisting}

\noindent\textbf{Pure-Rust Networking with Rustls}:
Instead of linking against native OpenSSL via C bindings, Fusion uses \texttt{rustls} with native root certificates. Network requests are performed asynchronously via Tokio and pure-Rust TLS, achieving complete immunity to Android Bionic libc linking issues while reducing resident memory to $<20$~MB.

\subsection{Bitwise Invariant Prefix Caching}
Modern high-performance LLM gateways (e.g., DeepSeek, vLLM Chunked Prefill, Anthropic Claude) feature hardware-level Key-Value (KV) cache reuse. In KV cache engines:
\begin{equation}
\text{KV Cache Hit}(L) \iff \forall i \in [0, L], \quad \text{token}_i^{(t)} \equiv \text{token}_i^{(t-1)}
\end{equation}
Any modified byte at index $k$ invalidates the cached KV tensors for all subsequent positions $i \ge k$.

In conventional conversational agents, because new messages, tool invocations, and thinking tokens are injected into the middle of the conversation history, the KV cache pointer is frequently invalidated on subsequent turns.

In Fusion, the leading prefix $P$ (which defines tool contracts, persona directives, and workspace specifications) is enforced to be **100\% bitwise invariant** from token index 0. Dynamic execution state $\Sigma_t$ and the immediate observation $O_t$ are appended strictly at the tail of the message payload. Consequently, the KV cache for the entire prompt prefix $P$ remains a guaranteed 100\% warm hit across turn 2 through turn 100+, yielding:
\begin{itemize}
    \item 80\%--90\% pricing discount on cached input tokens.
    \item Drastic reduction in Time-to-First-Token (TTFT) on mobile networks from $>3,500$~ms down to $<250$~ms.
\end{itemize}

\subsection{In-Browser WebAssembly \& Virtual Filesystem (VFS)}
To support zero-setup browser execution without containers or host OS access, Fusion compiles directly to WebAssembly (\texttt{wasm32-unknown-unknown}). 

\noindent\textbf{Dual-Tier Virtual Filesystem}:
Inside the browser tab, Fusion provides an in-memory VFS backed by IndexedDB persistence. All core filesystem operations (\texttt{read}, \texttt{write}, \texttt{edit}, \texttt{glob}, and \texttt{grep}) execute deterministically within browser memory. 

\noindent\textbf{Agent Client Protocol (ACP) Server}:
Fusion implements a complete JSON-RPC 2.0 ACP protocol server (\texttt{fusion --acp}) that runs identically across native Linux/Termux stdio and WebAssembly worker threads. When deployed inside cloud sandboxes (e.g., Cloudflare Containers, E2B microVMs), the headless agent runs in isolated Linux containers while streaming structured events to lightweight web or mobile frontends.
